% !TEX program = xelatex
% Шаблон научной статьи НИТУ «МИСиС»
% Заполните блоки ниже и скомпилируйте: xelatex article.tex

\documentclass{misisarticle}

% --- Классификационные коды (оставьте пустым {}, если не требуются) ---
\grnti{47.01.07, 47.01.09}
\asjc{1707, 2207}
\oecd{1020, 1070}

% --- Название ---
\title{Физико-информированные нейронные сети}

% --- Авторы: \misisauthor{ФИО}{аффилиация}{город, страна}{email} ---
\misisauthor{А.~О.~Виноградова}{кафедра инженерной кибернетики НИТУ «МИСиС»}{Москва, Россия}{m2312752@edu.misis.ru}
\misisauthor{Л.~А.~Артемьева}{кафедра инженерной кибернетики НИТУ «МИСиС»}{Москва, Россия}{supervisor@misis.ru}

% --- Аннотация и ключевые слова ---
\abstract{%
  Краткое изложение работы (5--10 предложений): постановка задачи, метод, основные результаты и выводы.
}

\keywords{%
  ключевое слово 1, ключевое слово 2, ключевое слово 3
}

\begin{document}

\maketitle

% ============================================================
\section{Введение}

Математическое моделирование процессов теплопереноса, диффузии, механики сплошных сред и других физических явлений во многих случаях сводится к решению краевых и начально-краевых задач для дифференциальных уравнений в частных производных (ДУЧП). Традиционные численные подходы, включая методы конечных разностей и конечных элементов, основаны на дискретизации расчетной области. Их применение хорошо изучено, однако построение сетки и увеличение числа степеней свободы могут осложнять вычисления в задачах со сложной геометрией, неоднородными коэффициентами и высокой размерностью.

Одним из подходов к решению ДУЧП без традиционной сеточной аппроксимации искомой функции стали физически информированные нейронные сети (Physics-Informed Neural Networks, PINN). В данном подходе решение представляется нейронной сетью, а уравнение и дополнительные условия включаются в функционал потерь посредством невязок, вычисляемых в наборе точек расчетной области. Метод позволяет объединять физические закономерности и наблюдаемые данные в единой процедуре оптимизации~\cite{raissi2019}. Вместе с тем совместное обучение всех параметров сети обычно приводит к нелинейной задаче оптимизации. Ее решение может быть чувствительно к выбору весов слагаемых функционала, распределению точек коллокации и обусловленности возникающей задачи~\cite{rathore2024}.

Альтернативный принцип обучения реализуется в экстремальных обучающих машинах (Extreme Learning Machines, ELM): параметры скрытых слоев задаются заранее, а оптимизации подлежит только линейный выходной слой~\cite{huang2006}. При использовании невязок дифференциального уравнения такой подход приводит к физически информированным ELM (PIELM)~\cite{dwivedi2020}. Дальнейшее развитие этой идеи связано с разбиением расчетной области на подобласти и представлением искомого решения совокупностью локальных сетей. В методе local Extreme Learning Machines (locELM) локальные представления объединяются условиями сопряжения, а их выходные веса определяются из задачи наименьших квадратов~\cite{dongli2021}.

В данной работе рассматривается связь глобальной PINN-аппроксимации с кусочным нейросетевым представлением на основе locELM. Основной интерес представляет построение гибридного алгоритма, в котором предварительно обученная глобальная PINN используется для выбора и инициализации скрытых признаков локальных сетей, после чего вычисляются новые выходные веса с учетом дифференциального уравнения, краевых условий и условий сопряжения. Такая схема предлагается как способ объединить адаптивное формирование признаков в PINN с линейным по выходным весам этапом обучения locELM. Ее вычислительные преимущества и ограничения должны устанавливаться в отдельных численных экспериментах.

\section{Анализ существующих работ}

В работе M.~Raissi, P.~Perdikaris и G.~E.~Karniadakis~\cite{raissi2019} предложена общая постановка PINN для прямых и обратных задач, в которой дифференциальные ограничения встраиваются в процедуру обучения. Для приближения искомой функции используется дифференцируемая нейронная сеть; производные по пространственным и временным координатам вычисляются автоматическим дифференцированием. Близкий по цели метод Deep Galerkin Method (DGM) применяет обучение по значениям невязок в случайно выбираемых точках области~\cite{sirignano2018}. Преимущество этих подходов заключается в гибкости выбора аппроксимирующей функции, однако обучение глубоких сетей может требовать большого числа итераций.

Особенности оптимизации PINN проанализированы в работе P.~Rathore и соавторов~\cite{rathore2024}. Авторы связывают трудности обучения, в частности, с плохой обусловленностью функционала, возникающей из-за действия дифференциального оператора. Рассмотрены комбинации методов оптимизации первого и второго порядка. Теоретическую основу для изучения поведения нейронных сетей при градиентном обучении также дают результаты A.~Jacot, F.~Gabriel и C.~Hongler о нейронном касательном ядре (Neural Tangent Kernel, NTK)~\cite{jacot2018}. Эти результаты не означают, что произвольная PINN эквивалентна линейной модели; вместе с тем они мотивируют исследование представлений, в которых часть параметров фиксируется, а заключительный этап сводится к линейной задаче.

В работе G.-B.~Huang, Q.-Y.~Zhu и C.-K.~Siew~\cite{huang2006} описан подход ELM для сетей с одним скрытым слоем: скрытые веса и смещения выбираются независимо от последующей оптимизации, а веса выходного слоя вычисляются линейными методами. V.~Dwivedi и B.~Srinivasan~\cite{dwivedi2020} перенесли эту конструкцию на решение ДУЧП, сформировав систему наименьших квадратов из значений невязки уравнения и дополнительных условий. Отмеченное авторами ограничение глобального PIELM на областях со сложным поведением решения послужило основанием для распределенного варианта метода.

S.~Dong и Z.~Li~\cite{dongli2021} разработали метод locELM, сочетающий локальные нейронные сети, декомпозицию области и условия сопряжения на внутренних границах. В их постановке параметры скрытых слоев фиксированы, а выходные коэффициенты подбираются из линейной либо нелинейной задачи наименьших квадратов в зависимости от вида дифференциального уравнения. Авторами рассмотрены стационарные и нестационарные задачи, а также блочная схема продвижения по времени. Таким образом, locELM позволяет увеличить гибкость локального представления, не переходя к одновременному градиентному обучению всех параметров сети.

Декомпозиция области используется и в других нейросетевых методах. В conservative PINNs (cPINNs) отдельные сети на подобластях согласуются с учетом законов сохранения~\cite{jagtap2020cpinn}; в extended PINNs (XPINNs) разбиение распространяется на пространственно-временную область и различные классы ДУЧП~\cite{jagtap2020xpinn}. Существенное различие между этими подходами и locELM состоит в организации оптимизации: в cPINNs и XPINNs параметры локальных PINN обучаются градиентными методами, тогда как в линейной постановке locELM при фиксированных скрытых слоях определяется только совокупность выходных коэффициентов.

Таким образом, опубликованные исследования обосновывают два независимых приема: адаптивное обучение нелинейных признаков в PINN и решение линейной задачи относительно выходных весов в locELM. В настоящей работе предлагается объединить эти приемы посредством переноса признаков предварительно обученной глобальной PINN в локальные ELM. В отличие от обычной случайной инициализации скрытых слоев locELM, этот вариант использует признаки, ранее настроенные на рассматриваемое дифференциальное уравнение. Данное сочетание рассматривается как \textit{предлагаемая модификация}, а не как уже экспериментально подтвержденное улучшение.

% ============================================================
\section{Постановка задачи}

\subsection{Исходная краевая задача}

Пусть $\Omega\subset\mathbb{R}^{2}$ --- ограниченная область с границей $\partial\Omega$. Требуется найти скалярную функцию $u:\overline{\Omega}\to\mathbb{R}$, удовлетворяющую линейному дифференциальному уравнению второго порядка и граничному условию Дирихле:
\begin{equation}
\begin{aligned}
\mathcal{L}u&=f, && \mathbf{x}\in\Omega,\\
 u&=g, && \mathbf{x}\in\partial\Omega,
\end{aligned}
\label{eq:generic}
\end{equation}
где $\mathbf{x}=(x,y)^{\mathsf T}$, $f$ --- заданная правая часть, $g$ --- граничная функция. Для конкретизации математического аппарата рассматривается эллиптический оператор
\begin{equation}
\mathcal{L}u=-\nabla\!\cdot\bigl(\kappa(\mathbf{x})\nabla u\bigr)
 +q(\mathbf{x})u,
\label{eq:operator}
\end{equation}
где $\kappa(\mathbf{x})\geq\kappa_{0}>0$ --- коэффициент диффузии, $q(\mathbf{x})\geq0$ --- заданный коэффициент. Для классической поточечной постановки предполагаются достаточная гладкость коэффициентов и решения; в частности, на каждой подобласти требуется возможность вычислять производные, входящие в~\eqref{eq:operator}.

Частным случаем задачи~\eqref{eq:generic} является уравнение Пуассона $-\Delta u=f$ при $\kappa\equiv1$ и $q\equiv0$. В качестве воспроизводимой модельной задачи для будущего численного исследования можно использовать $\Omega=(0,1)^2$, $g=0$ и
\begin{equation}
\begin{aligned}
 f(x,y)&=2\pi^2\sin(\pi x)\sin(\pi y),\\
 u_{\mathrm{ex}}(x,y)&=\sin(\pi x)\sin(\pi y).
\end{aligned}
\label{eq:poisson}
\end{equation}
Здесь $u_{\mathrm{ex}}$ является точным решением и может служить эталоном при последующей оценке погрешности. В настоящем разделе численные значения ошибок не приводятся.

\subsection{Глобальная нейросетевая аппроксимация}

В классическом подходе PINN искомое решение~\eqref{eq:generic} заменяется дифференцируемой функцией $u_{\theta}(\mathbf{x})$, задаваемой полносвязной нейронной сетью. Если последний слой линеен, ее можно представить в виде
\begin{equation}
 u_{\theta}(\mathbf{x})=a_{0}+\sum_{j=1}^{m}
 a_{j}\,\phi_{j}(\mathbf{x};\eta),
\label{eq:globalnet}
\end{equation}
где $\phi_j$ --- выходы последнего скрытого слоя, $\eta$ объединяет веса и смещения скрытых слоев, $a_0,a_j$ --- параметры выходного слоя, $\theta=(\eta,a_0,a_1,\ldots,a_m)$. При использовании гладкой функции активации, например $\tanh$, производные второго порядка могут быть вычислены автоматическим дифференцированием.

Введем наборы точек коллокации $X_f=\{\mathbf{x}_r^f\}_{r=1}^{N_f}\subset\Omega$ и $X_g=\{\mathbf{x}_r^g\}_{r=1}^{N_g}\subset\partial\Omega$. Тогда обучение PINN сводится к минимизации
\begin{equation}
\begin{aligned}
 J_{\mathrm{PINN}}(\theta)
 &=\frac{\lambda_f}{N_f}\sum_{r=1}^{N_f}
 \bigl|\mathcal{L}u_\theta(\mathbf{x}_r^f)-f_r\bigr|^2\\
 &\quad+\frac{\lambda_g}{N_g}\sum_{r=1}^{N_g}
 \bigl|u_\theta(\mathbf{x}_r^g)-g_r\bigr|^2,
\end{aligned}
\label{eq:pinnloss}
\end{equation}
где $\lambda_f,\lambda_g>0$ определяют относительную значимость внутренних и граничных невязок. В общем случае~\eqref{eq:pinnloss} нелинеен по $\eta$, поэтому для обучения используются итерационные методы оптимизации.

\subsection{Переход от глобальной к кусочной сети}

Разобьем расчетную область на $S$ непересекающихся по внутренности подобластей $\Omega_s$:
\begin{equation}
 \overline{\Omega}=\bigcup_{s=1}^{S}\overline{\Omega}_s,
 \qquad \Omega_i\cap\Omega_j=\varnothing\quad(i\ne j).
\label{eq:domain}
\end{equation}
На каждой подобласти введем собственную аппроксимацию
\begin{equation}
 u_h(\mathbf{x})\big|_{\Omega_s}=u_s(\mathbf{x})=
 \sum_{j=0}^{m_s}\beta_{sj}\phi_{sj}(\mathbf{x}),
\label{eq:localnet}
\end{equation}
где $\phi_{s0}\equiv1$, а остальные $\phi_{sj}$ --- фиксированные выходы скрытой части локальной сети. В отличие от глобальной PINN, неизвестными на локальном этапе являются лишь коэффициенты $\beta_{sj}$. Количество и вид базисных функций допускается выбирать отдельно для каждой подобласти.

Обозначим через $\Gamma_{ij}=\partial\Omega_i\cap\partial\Omega_j$ общую внутреннюю границу соседних подобластей, а через $\mathbf n_{ij}$ --- нормаль, направленную из $\Omega_i$ в $\Omega_j$. Для непрерывного решения задачи второго порядка в дивергентной форме естественными условиями сопряжения служат
\begin{equation}
\begin{aligned}
 u_i-u_j&=0, && \mathbf{x}\in\Gamma_{ij},\\
 \kappa_i\nabla u_i\cdot\mathbf n_{ij}
 -\kappa_j\nabla u_j\cdot\mathbf n_{ij}&=0,
 && \mathbf{x}\in\Gamma_{ij},
\end{aligned}
\label{eq:interfaces}
\end{equation}
где $\kappa_s=\kappa|_{\Omega_s}$; следовательно, на стыках согласуются значение решения и нормальный поток. При постоянном одинаковом $\kappa$ второе равенство выражает непрерывность нормальной производной. Если коэффициент диффузии имеет скачок, требование равенства самих нормальных производных, вообще говоря, некорректно: необходимо обеспечивать именно непрерывность потока.

Требуется построить кусочную функцию~\eqref{eq:localnet}, минимизирующую невязки~\eqref{eq:generic} внутри подобластей и на внешней границе, одновременно обеспечивая условия~\eqref{eq:interfaces} в выбранных интерфейсных точках. Дополнительно исследуется, как замена случайных скрытых весов локальных ELM на параметры, перенесенные из глобальной PINN, влияет на качество аппроксимации и обусловленность линейной системы.

\section{Описание методов}

\subsection{Базовый метод locELM}

В исходном подходе locELM~\cite{dongli2021} область разбивается на подобласти, каждой из которых соответствует отдельная нейронная сеть. Для сети с одним скрытым слоем фиксированные базисные функции могут иметь вид
\begin{equation}
 \phi_{sj}(\mathbf{x})=
 \sigma\bigl(\mathbf w_{sj}^{\mathsf T}
 \boldsymbol\xi_s(\mathbf{x})+b_{sj}\bigr),
\label{eq:features}
\end{equation}
где $\sigma$ --- гладкая функция активации, $\mathbf w_{sj}$ и $b_{sj}$ --- предварительно выбранные и затем неизменные параметры, $\boldsymbol\xi_s$ --- нормированные координаты внутри $\Omega_s$. Для нескольких скрытых слоев $\phi_{sj}$ понимается как $j$-й выход последнего из них. Фиксация скрытых параметров сохраняет линейность представления~\eqref{eq:localnet} относительно $\beta_{sj}$.

Внутри каждой подобласти выбираются точки $X_f^{(s)}$, на внешних границах --- точки $X_g^{(s)}$, а на общих границах соседних подобластей --- точки $X_{ij}^{\Gamma}$. В указанных точках записываются уравнение, граничное условие и оба условия сопряжения. Поскольку оператор~\eqref{eq:operator} линейный,
\begin{equation}
 \mathcal{L}u_s(\mathbf{x})=
 \sum_{j=0}^{m_s}\beta_{sj}
 \mathcal{L}\phi_{sj}(\mathbf{x}).
\label{eq:linearoperator}
\end{equation}
Поэтому все перечисленные уравнения линейны по искомым выходным весам. Для нелинейных дифференциальных операторов это свойство в общем случае утрачивается и требуется нелинейная задача наименьших квадратов~\cite{dongli2021}; далее рассматривается именно линейный эллиптический случай.

\subsection{Условия сопряжения и система наименьших квадратов}

Соберем выходные параметры всех подобластей в один вектор
\begin{equation}
 \boldsymbol\beta=
 (\boldsymbol\beta_1^{\mathsf T},\ldots,
 \boldsymbol\beta_S^{\mathsf T})^{\mathsf T},
 \quad \boldsymbol\beta_s=(\beta_{s0},\ldots,\beta_{sm_s})^{\mathsf T}.
\label{eq:coeffs}
\end{equation}
Запишем матрицы внутренних невязок $A_f$, граничных условий $A_g$ и условий сопряжения $A_0,A_1$. Например, строка $A_f$, относящаяся к точке $\mathbf{x}_r\in\Omega_s$, содержит в блоке $s$ значения $\mathcal{L}\phi_{sj}(\mathbf{x}_r)$ и нули в остальных блоках. Аналогично строка $A_g$ содержит значения $\phi_{sj}$ на внешней границе.

Для точки $\mathbf z_r\in\Gamma_{ij}$ строка $A_0$, соответствующая непрерывности решения, имеет коэффициенты $\phi_{i\ell}(\mathbf z_r)$ в блоке $i$ и $-\phi_{j\ell}(\mathbf z_r)$ в блоке $j$. Строка $A_1$ для согласования потока содержит соответственно
\begin{equation}
\begin{aligned}
 &\kappa_i(\mathbf z_r)\,
 \nabla\phi_{i\ell}(\mathbf z_r)\!\cdot\!\mathbf n_{ij},\\
 &-\kappa_j(\mathbf z_r)\,
 \nabla\phi_{j\ell}(\mathbf z_r)\!\cdot\!\mathbf n_{ij}.
\end{aligned}
\label{eq:interfaceblocks}
\end{equation}
Все прочие блоки этих строк равны нулю. Векторы $\mathbf f$ и $\mathbf g$ состоят из значений правой части и заданных граничных данных в соответствующих точках.

Пусть $N_f$, $N_g$, $N_0$, $N_1$ --- количества строк соответствующих групп ограничений. Взвешенный функционал имеет вид
\begin{equation}
\begin{aligned}
 J_{\mathrm{loc}}(\boldsymbol\beta)=
 &\frac{\lambda_f}{N_f}\|A_f\boldsymbol\beta-\mathbf f\|_2^2\\
 +&\frac{\lambda_g}{N_g}\|A_g\boldsymbol\beta-\mathbf g\|_2^2\\
 +&\frac{\lambda_0}{N_0}\|A_0\boldsymbol\beta\|_2^2\\
 +&\frac{\lambda_1}{N_1}\|A_1\boldsymbol\beta\|_2^2.
\end{aligned}
\label{eq:locloss}
\end{equation}
При отсутствии какой-либо группы ограничений соответствующее слагаемое исключается. Вводя диагональное взвешивание строк, задачу можно записать в компактной форме
\begin{equation}
 \boldsymbol\beta^{\ast}=
 \operatorname*{arg\,min}_{\boldsymbol\beta}
 \|A\boldsymbol\beta-\mathbf b\|_2^2,
\label{eq:ls}
\end{equation}
где
\begin{equation}
 A=\begin{bmatrix}
 \sqrt{\lambda_f/N_f}\,A_f\\
 \sqrt{\lambda_g/N_g}\,A_g\\
 \sqrt{\lambda_0/N_0}\,A_0\\
 \sqrt{\lambda_1/N_1}\,A_1
 \end{bmatrix},\quad
 \mathbf b=\begin{bmatrix}
 \sqrt{\lambda_f/N_f}\,\mathbf f\\
 \sqrt{\lambda_g/N_g}\,\mathbf g\\
 \mathbf 0\\\mathbf 0
 \end{bmatrix}.
\label{eq:stacked}
\end{equation}
Если $A$ имеет полный столбцовый ранг, решение~\eqref{eq:ls} единственно; в общем случае минимальное по евклидовой норме решение задается псевдообратной матрицей $\boldsymbol\beta^{\ast}=A^{+}\mathbf b$. На практике предпочтительно применять QR-разложение или сингулярное разложение (SVD), избегая явного формирования нормальных уравнений $A^{\mathsf T}A\boldsymbol\beta=A^{\mathsf T}\mathbf b$, которые могут усиливать проблемы обусловленности.

\subsection{Предлагаемая гибридная инициализация PINN--locELM}

Предлагаемая модификация состоит из двух последовательных этапов. На первом этапе обучается одна глобальная PINN в соответствии с~\eqref{eq:pinnloss}; в результате получаются скрытые параметры $\eta^{\ast}$ и выходные параметры $a^{\ast}=(a_0^{\ast},\ldots,a_m^{\ast})^{\mathsf T}$. На втором этапе область разделяется на подобласти, а в локальные сети переносятся веса и смещения скрытых слоев предварительно обученной модели:
\begin{equation}
 \phi_{sj}(\mathbf{x})=
 \phi_j(\mathbf{x};\eta^{\ast}),
 \qquad \mathbf{x}\in\Omega_s.
\label{eq:transfer}
\end{equation}
Веса скрытых слоев на втором этапе остаются фиксированными; оптимизируются только выходные коэффициенты $\boldsymbol\beta_s$. Таким образом, после предварительного нелинейного обучения глобальной сети локальная задача вновь сводится к линейному методу наименьших квадратов.

При одинаковом наборе перенесенных признаков и инициализации $\boldsymbol\beta_s^{(0)}=a^{\ast}$ для всех $s$ начальная кусочная аппроксимация в точности совпадает с глобальной сетью на каждой подобласти:
\begin{equation}
 u_s^{(0)}(\mathbf{x})=u_{\theta^{\ast}}(\mathbf{x})
 \quad(\mathbf{x}\in\Omega_s).
\label{eq:initialidentity}
\end{equation}
Следовательно, при достаточной гладкости глобальной сети ее начальное кусочное представление согласовано на интерфейсах по значению и по потоку при едином непрерывном коэффициенте $\kappa$. Это свойство сохраняется до изменения локальных выходных весов; поэтому в последующем решении~\eqref{eq:ls} интерфейсные ограничения должны быть включены явно.

Для ограничения отклонения от предварительно обученного представления предлагается исследовать регуляризованную постановку
\begin{equation}
 \boldsymbol\beta^{\ast}=\operatorname*{arg\,min}_{\boldsymbol\beta}
 \left\{\|A\boldsymbol\beta-\mathbf b\|_2^2
 +\mu\|\boldsymbol\beta-\boldsymbol\beta^{(0)}\|_2^2\right\},
\label{eq:regularized}
\end{equation}
где $\boldsymbol\beta^{(0)}$ образован повторением $a^{\ast}$, $\mu\geq0$ --- параметр регуляризации. Для $\mu>0$ эта задача эквивалентна расширенной системе наименьших квадратов
\begin{equation}
 \min_{\boldsymbol\beta}\left\|
 \begin{bmatrix}A\\\sqrt\mu\,I\end{bmatrix}
 \boldsymbol\beta-
 \begin{bmatrix}\mathbf b\\\sqrt\mu\,\boldsymbol\beta^{(0)}\end{bmatrix}
 \right\|_2^2.
\label{eq:augmented}
\end{equation}
Регуляризация обеспечивает единственность решения при $\mu>0$ и ограниченном числе параметров, но может создавать смещение и ухудшать выполнение дифференциальных ограничений. Поэтому значение $\mu$ должно подбираться по независимым контрольным точкам, а результат необходимо сопоставлять с вариантом $\mu=0$.

Для проверки гипотезы о полезности предварительного обучения рассматриваются два сопоставимых варианта локальных сетей: (а)~стандартный locELM со случайно заданными скрытыми параметрами и (б)~PINN--locELM с перенесенными скрытыми параметрами. Для корректного сравнения необходимо фиксировать число подобластей, число выходных коэффициентов, точки коллокации, способ решения линейной системы и бюджет вычислений. Перенос признаков не гарантирует повышения точности: после глобального обучения признаки могут оказаться линейно зависимыми или недостаточно выразительными для отдельных подобластей. Эти эффекты подлежат численной проверке.

\subsection{Возможные дополнения и критерии проверки}

При использовании локальных координат $\boldsymbol\xi_s$ важно не менять перенесенную функцию без учета преобразования аргумента. Например, если $\mathbf{x}=\mathbf c_s+D_s\boldsymbol\xi_s$ и первый скрытый слой глобальной сети имеет вид $\sigma(W\mathbf{x}+\mathbf b)$, то эквивалентный локальный слой записывается как
\begin{equation}
 \sigma\!\left((WD_s)\boldsymbol\xi_s
 +(W\mathbf c_s+\mathbf b)\right).
\label{eq:coordtrans}
\end{equation}
Это позволяет нормировать координаты подобластей, сохраняя точное совпадение начальной локальной и глобальной моделей. При многослойной сети указанное преобразование затрагивает первый аффинный слой; остальные скрытые слои могут быть перенесены без изменения.

В качестве дополнительного приема предлагается исследовать масштабирование блоков~\eqref{eq:stacked} и, при необходимости, отбор перенесенных признаков с помощью рангообнаруживающего QR-разложения. Такое сокращение числа локальных признаков способно уменьшить размер системы, однако нарушает тождество~\eqref{eq:initialidentity}; после отбора коэффициенты следует переопределить и заново обеспечить сопряжение. Дополнительное требование согласования касательных производных на интерфейсе может рассматриваться как дискретная стабилизация при гладком решении, но не заменяет физически необходимого условия непрерывности нормального потока.

Эффективность подхода предполагается оценивать по относительной ошибке $L^2$ на независимой контрольной выборке, максимальной невязке условий сопряжения, числу обусловленности системы и времени вычислений. Для модельной задачи~\eqref{eq:poisson} относительная ошибка может быть определена как
\begin{equation}
 \varepsilon_{\mathrm{rel}}=
 \frac{\left(\sum_{r=1}^{N_{\mathrm{test}}}
 |u_h(\mathbf z_r)-u_{\mathrm{ex}}(\mathbf z_r)|^2\right)^{1/2}}
 {\left(\sum_{r=1}^{N_{\mathrm{test}}}
 |u_{\mathrm{ex}}(\mathbf z_r)|^2\right)^{1/2}}.
\label{eq:relativeerror}
\end{equation}
Предполагается отдельно сравнивать время предварительного обучения PINN, время локального решения и их сумму: ускорение только второго этапа не означает ускорения полного алгоритма. Варьирование граничных данных, коэффициента $\kappa$ и числа подобластей позволит проверить устойчивость метода на классе задач, а не на одном частном примере.

% ============================================================
\section{Заключение}

Кратко сформулируйте основные выводы работы.

% ============================================================
% Список литературы — вариант 1: вручную
begin{thebibliography}{9}

\bibitem{raissi2019}
Raissi M., Perdikaris P., Karniadakis G.~E. Physics-informed neural networks: A deep learning framework for solving forward and inverse problems involving nonlinear partial differential equations. \textit{Journal of Computational Physics}, 2019, vol.~378, pp.~686--707. DOI:~\url{10.1016/j.jcp.2018.10.045}.

\bibitem{sirignano2018}
Sirignano J., Spiliopoulos K. DGM: A deep learning algorithm for solving partial differential equations. \textit{Journal of Computational Physics}, 2018, vol.~375, pp.~1339--1364. DOI:~\url{10.1016/j.jcp.2018.08.029}.

\bibitem{rathore2024}
Rathore P., Lei W., Frangella Z., Lu L., Udell M. Challenges in Training PINNs: A Loss Landscape Perspective. \textit{Proceedings of Machine Learning Research}, 2024, vol.~235, pp.~42159--42191. URL:~\url{https://proceedings.mlr.press/v235/rathore24a.html}.

\bibitem{jacot2018}
Jacot A., Gabriel F., Hongler C. Neural Tangent Kernel: Convergence and Generalization in Neural Networks. \textit{Advances in Neural Information Processing Systems}, 2018, vol.~31. URL:~\url{https://arxiv.org/abs/1806.07572}.

\bibitem{huang2006}
Huang G.-B., Zhu Q.-Y., Siew C.-K. Extreme learning machine: Theory and applications. \textit{Neurocomputing}, 2006, vol.~70, no.~1--3, pp.~489--501. DOI:~\url{10.1016/j.neucom.2005.12.126}.

\bibitem{dwivedi2020}
Dwivedi V., Srinivasan B. Physics Informed Extreme Learning Machine (PIELM)---A rapid method for the numerical solution of partial differential equations. \textit{Neurocomputing}, 2020, vol.~391, pp.~96--118. DOI:~\url{10.1016/j.neucom.2019.12.099}.

\bibitem{dongli2021}
Dong S., Li Z. Local extreme learning machines and domain decomposition for solving linear and nonlinear partial differential equations. \textit{Computer Methods in Applied Mechanics and Engineering}, 2021, vol.~387, art.~114129. DOI:~\url{10.1016/j.cma.2021.114129}.

\bibitem{jagtap2020cpinn}
Jagtap A.~D., Kharazmi E., Karniadakis G.~E. Conservative physics-informed neural networks on discrete domains for conservation laws: Applications to forward and inverse problems. \textit{Computer Methods in Applied Mechanics and Engineering}, 2020, vol.~365, art.~113028. DOI:~\url{10.1016/j.cma.2020.113028}.

\bibitem{jagtap2020xpinn}
Jagtap A.~D., Karniadakis G.~E. Extended Physics-Informed Neural Networks (XPINNs): A Generalized Space-Time Domain Decomposition Based Deep Learning Framework for Nonlinear Partial Differential Equations. \textit{Communications in Computational Physics}, 2020, vol.~28, no.~5, pp.~2002--2041. DOI:~\url{10.4208/cicp.OA-2020-0164}.

\end{thebibliography}
\end{document}

